\chapter{Scrum Charter}

\section{Product goal}

Deliver an AI-native customer support platform that handles complete text conversations first, provides grounded answers and safe tool use, supports human handoff, and then expands to real-time voice and additional channels.

\section{Scrum setup}

\begin{itemize}
\item \textbf{Product Owner:} accountable for product value, Product Goal, and Product Backlog ordering.
\item \textbf{Scrum Master:} coaches Scrum, removes impediments, and helps the team improve.
\item \textbf{Developers:} cross-functional team that designs, builds, tests, integrates, secures, and operates a usable Increment.
\item \textbf{Stakeholders:} support teams, organization administrators, customers, security/privacy, and business-system owners.
\end{itemize}
Roles can be held by a small team, but accountabilities must remain clear.

\mmdfigL{07-scrum-release-risks/01-7-1-scrum-roles-and-accountabilities}{Scrum roles and accountabilities}{fig:ease-risks-01-7-1-scrum-roles-and-accountabilities-3}


\section{Product Goal}

A customer can contact an organization through website chat, receive accurate knowledge-grounded answers, complete at least one authorized read-only business lookup, and reach a human without losing conversation context. The architecture must be ready to add voice without duplicating business reasoning.

\section{Proposed release goals}

\begin{itemize}
\item \textbf{Release 1 --- Text support MVP:} tenant foundation, website chat, conversation lifecycle, RAG, one read-only tool, handoff, basic analytics and audit.
\item \textbf{Release 2 --- Voice MVP:} real-time browser voice, speech processing, transcript synchronization, interruption handling, voice handoff and voice observability.
\item \textbf{Release 3 --- Operations expansion:} additional channels, workflow automation, routing/SLA, connectors, evaluation, and commercial usage controls.
\end{itemize}
\section{Sprint approach}

\begin{itemize}
\item Suggested Sprint length: 2 weeks.
\item Each Sprint has one Sprint Goal and a potentially releasable Increment.
\item Sprint Planning selects work based on Product Goal, ordered backlog, capacity, and Definition of Done.
\item Daily Scrum is for Developers to inspect progress toward the Sprint Goal and adapt the plan.
\item Sprint Review inspects the Increment with stakeholders and adapts the Product Backlog.
\item Sprint Retrospective improves team effectiveness and quality.
\item Refinement is ongoing; it is not a formal Scrum event.
\end{itemize}
\mmdfigL{07-scrum-release-risks/02-7-2-scrum-events-within-a-sprint}{Scrum events within a Sprint}{fig:-release-risks-02-7-2-scrum-events-within-a-sprint-4}


\mmdfig{scrum-existing/S5-scrum-mermaid-artifacts/04-4-scrum-workflow}{Scrum workflow}{fig:ing-s5-scrum-mermaid-artifacts-04-4-scrum-workflow-5}


\section{Working assumptions}

\begin{itemize}
\item A small team is building a modular monolith initially.
\item Product design, infrastructure, security and automated tests are part of delivery---not separate downstream phases.
\item Story points are team-specific estimates, not time conversions.
\item Sprint sequence and estimates are provisional until capacity and dependencies are confirmed.
\end{itemize}

\section{Scrum Events, Artifacts and Working Agreements}

\subsection{Sprint Planning}

Inputs:

\begin{itemize}
\item Product Goal and ordered Product Backlog
\item latest Increment and feedback
\item team capacity, leave, and known operational work
\end{itemize}
Outputs:

\begin{itemize}
\item Sprint Goal
\item selected Product Backlog Items
\item initial delivery plan
\end{itemize}
\subsection{Daily Scrum}

\begin{itemize}
\item 15 minutes for Developers.
\item Inspect progress toward the Sprint Goal.
\item Adapt the plan for the next day.
\item Use follow-up conversations for detailed problem solving.
\end{itemize}
\subsection{Sprint Review}

\begin{itemize}
\item Demonstrate working software, not slides alone.
\item Inspect outcomes against the Sprint Goal.
\item Invite support staff/admins to test realistic scenarios.
\item Capture feedback and reorder backlog based on evidence.
\end{itemize}
\subsection{Sprint Retrospective}

Inspect:

\begin{itemize}
\item delivery flow and collaboration;
\item defect/rework causes;
\item CI/test speed;
\item architectural friction;
\item security/privacy and operational issues;
\item one or two actionable improvements for the next Sprint.
\end{itemize}
\subsection{Backlog refinement}

Ongoing activity to split, clarify, reorder, and estimate upcoming work. Refinement is not a formal Scrum event.

\subsection{Scrum artifacts and commitments}

\begin{itemize}
\item \textbf{Product Backlog / Product Goal:} ordered work and desired future state.
\item \textbf{Sprint Backlog / Sprint Goal:} selected work and plan to achieve the Sprint Goal.
\item \textbf{Increment / Definition of Done:} integrated usable result meeting quality requirements.
\end{itemize}
\subsection{Suggested working agreements}

\begin{itemize}
\item Every story has observable acceptance criteria.
\item Security, observability, and automated tests are included in feature work.
\item Merge small changes frequently.
\item No direct production secrets in source control.
\item Do not let model output directly bypass application permissions.
\item Review incidents and failed evaluations without blaming individuals.
\item Keep story points team-specific; never treat them as hours.
\end{itemize}
